\section{Valuation profiles and integral Fourier data}
Let $N=\prod_{i=1}^h p_i^{a_i}$ with distinct primes $p_i$ and positive $a_i$.  We allow $h=0$ when $N=1$.  Via the Chinese remainder theorem, write $\mathbb Z_N\cong\prod_i\mathbb Z_{p_i^{a_i}}$ and let
\[
 \Lambda_N=\prod_{i=1}^h\{0,\ldots,a_i\},\qquad
 v_i(x)=\min(a_i,v_{p_i}(x_i)),\quad v_{p_i}(0)=a_i.
\]
The order on $\Lambda_N$ is coordinatewise.  A function $f:\mathbb Z_N\to\mathbb Z$ is \emph{radial} if $f(x)=f(\vv(x))$ depends only on the capped valuation vector.  Its profile will be written $f(\rvec)$.  Put
\[
 H_{\rvec}=\prod_i p_i^{r_i}\mathbb Z_{p_i^{a_i}},\qquad
 W_{\rvec}=|H_{\rvec}|=\prod_i p_i^{a_i-r_i}.
\]
Thus $H_{\rvec}$ consists of all points with valuation vector $\vv\ge\rvec$.  Fourier transforms are unnormalized:
\[
 \widehat f(\chi)=\sum_{x\in\mathbb Z_N} f(x)\overline{\chi(x)},
 \qquad f(x)=N^{-1}\sum_{\chi}\widehat f(\chi)\chi(x).
\]
A character has \emph{conductor} $c(\s)=\prod_i p_i^{s_i}$ for a unique $\s\in\Lambda_N$.  A transform of a radial function is constant on each conductor class; we denote its value by $F(\s)$, or $F(c)$ if convenient.  There are exactly $\varphi(c(\s))$ characters in that class.

\begin{lemma}[Subgroup-union profiles]\label{ii:lem:upset}
A set $A\subseteq\mathbb Z_N$ is a union of additive subgroups if and only if its valuation profile is a $\{0,1\}$-valued increasing function on $\Lambda_N$.  The empty set and the full group are included.
\end{lemma}
\begin{proof}
Each subgroup $H_{\rvec}$ is the upper principal set $\{\vv:\vv\ge\rvec\}$ in valuation space.  Their unions are increasing and binary.  Conversely, every increasing subset of the finite poset $\Lambda_N$ is a union of the principal upper sets generated by its minimal elements.  The Chinese remainder theorem identifies these with the additive subgroups $H_{\rvec}$.\qedhere
\end{proof}

\begin{lemma}[Integral conductor inversion]\label{ii:lem:integrality}
An integer radial function $f$ admits a unique integral expansion
\begin{equation}\label{ii:eq:alpha}
 f=\sum_{\rvec\in\Lambda_N}\alpha_{\rvec}\one_{H_{\rvec}},\qquad
 \alpha_{\rvec}=\sum_{J\subseteq\{i:r_i>0\}}(-1)^{|J|}f(\rvec-\one_J).
\end{equation}
Its conductor coefficients are integers and satisfy
\begin{equation}\label{ii:eq:Fs}
 F(\s)=\sum_{\rvec\ge\s}W_{\rvec}\alpha_{\rvec},
 \qquad
 F(\s)-F(\s+\e_i)\equiv0\pmod{p_i^{a_i-s_i}}
 \quad(s_i<a_i).
\end{equation}
Conversely, an arbitrary integer array $G:\Lambda_N\to\mathbb Z$ is the conductor transform of an integer radial function if and only if it satisfies every adjacent congruence in \eqref{ii:eq:Fs}.  Its coefficients are explicitly
\begin{equation}\label{ii:eq:inversealpha}
 \alpha_{\rvec}=\frac{1}{W_{\rvec}}
 \sum_{J\subseteq\{i:r_i<a_i\}}(-1)^{|J|}G(\rvec+\one_J).
\end{equation}
\end{lemma}
\begin{proof}
M\"obius inversion on the product of chains yields \eqref{ii:eq:alpha}, since $f(\vv)=\sum_{\rvec\le\vv}\alpha_{\rvec}$.  The transform of $\one_{H_{\rvec}}$ is $W_{\rvec}$ on characters trivial on $H_{\rvec}$, i.e., with conductor vector $\s\le\rvec$, and zero otherwise.  This gives the first formula in \eqref{ii:eq:Fs}.  Subtracting adjacent entries retains only the terms with $r_i=s_i$, each divisible by $p_i^{a_i-s_i}$.

Conversely, repeated finite differences in the cumulative identity give \eqref{ii:eq:inversealpha}.  When $r_i<a_i$, group its numerator into differences along coordinate $i$; each difference is divisible by $p_i^{a_i-r_i}$.  This holds for every such $i$, and these prime-power divisors are pairwise coprime.  Therefore their product $W_{\rvec}$ divides the numerator, so each $\alpha_{\rvec}$ is integral.  The cumulative transform recovers $G$ exactly, and uniqueness follows from the inversion formula.  This also proves that there are no additional mixed-coordinate integrality conditions.\qedhere
\end{proof}

\begin{remark}[Connection to cyclic Burnside marks]\label{ii:rem:burnside}
Regard $\widehat{\mathbb Z_N}$ as a cyclic group.  For $H_{\rvec}\le\mathbb Z_N$, the annihilator $H_{\rvec}^{\perp}$ is a subgroup of the dual group, and the transitive dual-group set $\widehat{\mathbb Z_N}/H_{\rvec}^{\perp}$ has $|H_{\rvec}|=W_{\rvec}$ points.  Its mark at the subgroup generated by a character of conductor $c(\s)$ is $W_{\rvec}$ when $\s\le\rvec$, and zero otherwise.  These are precisely the columns in \eqref{ii:eq:Fs}.  Therefore the converse in Lemma~\ref{ii:lem:integrality} describes the cyclic Burnside mark lattice in these coordinates; its local congruences are a prime-primary presentation of the classical mark-image congruences \cite{tomDieck1979,BoltjePfeiffer2008}.  We make no novelty claim for this integrality characterization.
\end{remark}

\begin{remark}\label{ii:rem:fourier}
Every radial function is invariant under $x\mapsto-x$.  Its Fourier coefficients are therefore real.  For an integer radial function they are in fact integers by Lemma~\ref{ii:lem:integrality}.  The convolution identity gives $\widehat{C_f}(\chi)=|\widehat f(\chi)|^2$.  Thus two radial real functions have the same complete autocorrelation exactly when their conductor coefficient squares agree for every $\s$.  In Parseval's identity \emph{each conductor is repeated $\varphi(c(\s))$ times}.
\end{remark}

\section{Fourier sign graphs and exact radial reconstruction}
Fix a nonempty radial binary set $A$ with conductor transform $F$.  Define the \emph{sign graph} $\Gamma_A$ as follows:
\begin{itemize}[nosep]
 \item Its vertices are exactly $V_A=\{\s\in\Lambda_N:F(\s)\ne0\}$.
 \item Two vertices $\s$ and $\s+\e_i$ ($s_i<a_i$) are joined if
 \begin{equation}\label{ii:eq:edge}
 p_i^{a_i-s_i}\nmid 2F(\s).
 \end{equation}
\end{itemize}
The test in \eqref{ii:eq:edge} is symmetric in the endpoints by \eqref{ii:eq:Fs}.  A vertex is \emph{root-marked} if $\s=\boldsymbol0$ and \emph{pure-prime-marked} if its conductor is a positive power of a single prime.  Components containing either kind of marked vertex are called \emph{divisor-anchored}.  Define
\begin{align*}
 u_{\rm rad}(A)&=\#\{\text{components not containing }\boldsymbol0\},\\
 u_{\rm div}(A)&=\#\{\text{components containing no root or pure-prime mark}\}.
\end{align*}
The graph is formed from the \emph{actual} conductor coefficients of the specified $A$, not an abstract list of allowable signs.

\begin{theorem}[Exact radial binary sign classification]\label{ii:thm:radial}
Let $A$ be any nonempty radial binary subset of $\mathbb Z_N$.  Its radial binary homometric mates are in bijection with the subsets of the $u_{\rm rad}(A)$ non-root components of $\Gamma_A$.  The mate corresponding to a selection is obtained by replacing $F(\s)$ with $-F(\s)$ at every vertex in each selected component, preserving $F$ elsewhere, and applying Fourier inversion.  In particular the number of radial binary mates, including $A$, is exactly $2^{u_{\rm rad}(A)}$.
\end{theorem}
\begin{proof}
For any radial binary mate $b$, \eqref{ii:eq:corr} fixes both its Fourier magnitudes and its cardinality.  Hence its integer conductor transform $G$ satisfies $G(\s)=\epsilon_{\s}F(\s)$ for $\s\in V_A$, with $\epsilon_{\s}\in\{\pm1\}$, and is zero off $V_A$.  Since $G(\boldsymbol0)=|A|=F(\boldsymbol0)>0$, the root sign is positive.  At adjacent nonzero vertices, the two congruences \eqref{ii:eq:Fs} for $F$ and $G$ show that opposite signs would require $p_i^{a_i-s_i}\mid2F(\s)$.  Thus the signs agree along each edge, and a mate determines a choice of non-root components.

For the converse, flip any chosen non-root components and let $G$ be the resulting integer array.  We verify \eqref{ii:eq:Fs} for $G$ at every adjacent pair with modulus $m=p_i^{a_i-s_i}$.  If both entries are nonzero and the edge \eqref{ii:eq:edge} exists, their signs are the same.  If the edge does not exist, $m\mid2F(\s)$ and, by $F(\s)\equiv F(\s+\e_i)\pmod m$, also $m\mid2F(\s+\e_i)$; either sign choice preserves the congruence.  If exactly one entry is zero, the other entry was already divisible by $m$, so either sign is allowed.  Two zeros are automatic.  Lemma~\ref{ii:lem:integrality} therefore gives a unique \emph{integer radial} inverse $b$.

Because the root sign was not changed,
\begin{equation}\label{ii:eq:mass}
 \sum_x b(x)=G(\boldsymbol0)=|A|.
\end{equation}
By Parseval, with the correct conductor multiplicities,
\begin{equation}\label{ii:eq:energy}
 \sum_x b(x)^2=\frac1N\sum_{\s\in\Lambda_N}
 \varphi(c(\s))G(\s)^2
 =\frac1N\sum_{\s}\varphi(c(\s))F(\s)^2=|A|.
\end{equation}
Subtracting \eqref{ii:eq:mass} from \eqref{ii:eq:energy} gives $\sum_x b(x)(b(x)-1)=0$.  Every integer $z$ satisfies $z(z-1)\ge0$, with equality exactly for $z\in\{0,1\}$, even when negative integers are allowed.  Thus $b$ is binary.  Its squared Fourier transform equals that of $\one_A$ at every character, so it is homometric to $A$.

Two different component selections change at least one nonzero conductor coefficient, and therefore give different functions by Fourier uniqueness.  Necessity proved that every radial binary mate occurs.\qedhere
\end{proof}

\section{Odd-order uniqueness and trivial sign flips}

\begin{proposition}[Symmetric sets at odd modulus]\label{ii:prop:odd}
If $N$ is odd and two binary subsets $A,B\subseteq\mathbb Z_N$ satisfy $A=-A$, $B=-B$, and $C_A=C_B$, then $A=B$.  In particular $u_{\rm rad}(A)=0$ for every nonempty radial binary set of odd modulus.
\end{proposition}
\begin{proof}
For a symmetric binary indicator $f$, $C_f=f*f$.  In $\mathbb F_2[\mathbb Z_N]$ the off-diagonal terms cancel in pairs $(x,y)$ and $(y,x)$, so
\[
 C_f(t)\equiv\sum_{2x=t}f(x)^2\equiv f(t/2)\pmod2.
\]
Doubling is bijective because $N$ is odd.  Thus the autocorrelation determines $f$.  The final assertion follows from Theorem~\ref{ii:thm:radial}, since distinct non-root component flips would give distinct radial binary mates.
\end{proof}

\begin{proposition}[Translation and complementation components]\label{ii:prop:symmetries}
Let $A\subseteq\mathbb Z_N$ be nonempty and radial binary, with $u=u_{\rm rad}(A)$.  Let $d\mid N$ be its \emph{least positive period}, meaning $A+d=A$ and no smaller positive period exists.  Every radial binary mate has the same least period $d$.

\begin{enumerate}[label=(\roman*),nosep]
\item If $d$ is even, translation by $d/2$ maps each radial mate to a distinct radial mate and flips precisely the nonzero conductor coefficients for which $v_2(c)=v_2(d)$.  Those vertices form a nonempty union $T$ of non-root graph components.  If $d$ is odd, no nontrivial translation of a radial mate can be radial.
\item When $N$ is even and $|A|=N/2$, complementation maps each radial mate to a distinct radial mate and flips \emph{all} non-root sign components.  Outside the half-size case, complementation does not preserve the given autocorrelation.
\item Let $V\cong\mathbb F_2^u$ index the non-root sign components.  Write $\mathbf t\in V$ for the translation mask in (i), when $d$ is even, and $\mathbf1\in V$ for the full mask in (ii), when $|A|=N/2$.  The number of radial mates modulo translations, reflection, and (when homometric) complementation is exactly
\begin{equation}\label{ii:eq:orbitcount}
 2^{u-r},\qquad r=\dim_{\mathbb F_2}\left\langle\;\mathbf t\ (d\text{ even}),\quad\mathbf1\ (|A|=N/2)\;\right\rangle.
\end{equation}
Omitted masks are not generators; thus $0\le r\le2$.  Reflection fixes every radial set.
\end{enumerate}
\end{proposition}
\begin{proof}
The set of translation periods is a subgroup of $\mathbb Z_N$, generated by $d$.  Fourier inversion shows it is determined by the nonzero Fourier support: a shift is a period if and only if every character with nonzero coefficient takes the value $1$ there.  All radial homometric mates have the same support and therefore the same $d$.

If a translate $B+t$ of a radial mate $B$ is itself radial, then both $B$ and $B+t$ are invariant under negation.  Hence $B+t=B-t$, so $2t$ is a period, and $2t\equiv0\pmod d$.  Thus $t\equiv0\pmod d$ if $d$ is odd, and $t\equiv0,d/2\pmod d$ if $d$ is even.  Conversely, for even $d$, every unit $q$ modulo $N$ is odd; consequently $q(d/2)\equiv d/2\pmod d$.  Radiality and $d$-periodicity imply $B+d/2$ is radial.  Its Fourier coefficient at a character of conductor $c\mid d$ is multiplied by $\chi(d/2)$, equal to $-1$ exactly when $v_2(c)=v_2(d)$; all nonzero coefficients have $c\mid d$.  Since $d$ is the least period, some nonzero coefficient has this maximal $2$-adic conductor height.  Theorem~\ref{ii:thm:radial} then shows that the flip set $T$ must be a nonempty union of graph components, giving (i).  Independently of periodicity, no graph edge can cross from the penultimate to the top ambient $2$-adic conductor layer, because its modulus is $2$ and $2\mid2F(\s)$.  In particular, when $v_2(d)=v_2(N)$, the nonzero top ambient layer is a union of components and translation by $N/2$ flips precisely that layer.  If $d$ has smaller $2$-adic height, a smaller half-period, rather than $N/2$, is the relevant translation.

For a set of size $k$, $C_{A^c}(t)=N-2k+C_A(t)$ for all $t\ne0$, and $C_{A^c}(0)=N-k$.  It has the same autocorrelation exactly when $k=N/2$.  Then its root coefficient remains $k$ and every other Fourier coefficient is negated.  This proves (ii).  Translation and complementation act on the bit vectors of component signs by addition of the indicated masks.  Their subgroup has order $2^r$ and acts freely on the $2^u$ selections.  Reflection acts trivially.  The orbit count follows.
\end{proof}

\section{Recovering the subgroup-union structure}
The preceding theorem does \emph{not} guarantee that the radial binary inverse is increasing.  The next two lemmas provide the missing implication.

\begin{lemma}[Pure-prime spectral positivity]\label{ii:lem:pure}
If $A$ is a subgroup union and $p^a\Vert N$, then $F(p^j)\ge0$ for $0\le j\le a$.  Consequently every subgroup-union homometric mate has exactly the same coefficients at the root and every pure prime-power conductor.
\end{lemma}
\begin{proof}
Write $\mathbb Z_N\cong\mathbb Z_{p^a}\times Y$.  For $y\in Y$, let $u_r(y)\in\{0,1\}$ be the value of $\one_A$ when the $p$ coordinate has valuation $r$, $0\le r\le a$.  Increasingness gives $u_0\le\cdots\le u_a$.  With $u_{-1}=0$ put $\eta_r=u_r-u_{r-1}\ge0$.  Telescoping expresses the $p$-fiber as $\sum_{r=0}^a\eta_r(y)\one_{p^r\mathbb Z_{p^a}}$.  At a character of conductor $p^j$ trivial on $Y$, the subgroup sum equals $p^{a-r}$ for $r\ge j$, and zero otherwise.  Hence
\[
 F(p^j)=\sum_{r=j}^a p^{a-r}\sum_{y\in Y}\eta_r(y)\ge0.
\]
The Fourier magnitudes of homometric subgroup unions agree.  A real, nonnegative coefficient with specified magnitude is uniquely determined, including when it is zero.\qedhere
\end{proof}

\begin{lemma}[Marginal-autocorrelation defect]\label{ii:lem:defect}
Let $A$ be a subgroup union and $b$ a radial binary function homometric to $\one_A$ whose Fourier coefficients agree with those of $A$ at the root and every pure prime-power conductor.  Then $b$ is increasing in each valuation coordinate and hence is itself the indicator of a subgroup union.
\end{lemma}
\begin{proof}
Fix $p^a\Vert N$ and set $Y=\prod_{q\ne p}\mathbb Z_{q^{a_q}}$.  For $0\le r\le a$ let $b_r(y)$ denote the value of $b$ on the $p$-shell with exact capped valuation $r$, and define $M_r(b)=\sum_y b_r(y)$.  The marginal $m_b(z)=\sum_y b(z,y)$ on $\mathbb Z_{p^a}$ has Fourier transform given by the root and the $p^j$ conductor values.  All of these agree with the corresponding values for $A$, so ordinary one-dimensional Fourier inversion gives $m_b=m_A$, and hence
\begin{equation}\label{ii:eq:marginal}
 M_r(b)=M_r(A)\quad(0\le r\le a).
\end{equation}

For $r<a$ write $w_r=(p-1)p^{a-r-1}$, the number of elements in the shell of valuation $r$, and put $w_a=1$.  The same-shell counts below use $b_r(y)^2=b_r(y)$, since the shell values are binary.  Fix a shift $t=(t_p,0)$ with $v_p(t_p)=s<a$.  In the ordered-pair count $C_b(t)$, for each $y$:
\begin{enumerate}[label=(\roman*),nosep]
 \item points of valuation $r<s$ remain in shell $r$, giving $w_rb_r(y)$;
 \item points with both endpoints in shell $s$ number $(p-2)p^{a-s-1}$, giving $(p-2)p^{a-s-1}b_s(y)$ (this term is zero for $p=2$);
 \item a pair with one endpoint in shell $s$ and the other in shell $r>s$ gives $2w_rb_s(y)b_r(y)$.
\end{enumerate}
There are no further possibilities: adding $t_p$ to a lower-valuation point retains its valuation, while adding it to a higher-valuation point lands in shell $s$.  Thus
\begin{align}
 C_b(t)&=\sum_{r<s}w_rM_r(b)+(p-2)p^{a-s-1}M_s(b)
          +2\sum_{y\in Y}b_s(y)\sum_{r>s}w_rb_r(y).\label{ii:eq:shellcorr}
\end{align}
Since $\sum_{r>s}w_r=p^{a-s-1}$, rearranging gives the \emph{nonnegative defect identity}
\begin{equation}\label{ii:eq:defect}
 \underbrace{\left(\sum_{r<s}w_rM_r(b)+p^{a-s}M_s(b)-C_b(t)\right)}_{D_{p,s}(b)}
 =2\sum_{y\in Y}b_s(y)\sum_{r>s}w_r\bigl(1-b_r(y)\bigr)\ge0.
\end{equation}
For $A$ the right side is zero, because every fiber is increasing.  By \eqref{ii:eq:marginal} and homometry, $D_{p,s}(b)=D_{p,s}(A)=0$.  Every summand on the right of \eqref{ii:eq:defect} is nonnegative and every $w_r>0$; therefore $b_s(y)=1$ implies $b_r(y)=1$ for every $r>s$.  This applies to each prime-power coordinate, including the dyadic coordinate. Thus the profile is coordinatewise increasing, and Lemma~\ref{ii:lem:upset} applies.\qedhere
\end{proof}

\begin{theorem}[Complete classification among subgroup unions]\label{ii:thm:subgroups}
For every $N\ge1$ and every nonempty subgroup union $A\subseteq\mathbb Z_N$, its subgroup-union homometric mates are in bijection with the selections of components of $\Gamma_A$ containing no root or pure prime-power vertex.  In particular:
\begin{enumerate}[label=(\alph*),nosep]
\item every selection of these \emph{unanchored} components has a unique binary subgroup-union inverse with exactly the same complete autocorrelation;
\item every subgroup-union homometric mate arises from exactly one such selection;
\item $A$ is rigid among subgroup unions if and only if $u_{\rm div}(A)=0$, and the number of such mates including $A$ is $2^{u_{\rm div}(A)}$.
\end{enumerate}
\end{theorem}
\begin{proof}
Every subgroup union is radial binary, so by Theorem~\ref{ii:thm:radial} its radial homometric mates can only differ through sign changes on non-root components.  Lemma~\ref{ii:lem:pure} forces all nonzero pure-prime Fourier signs to remain positive.  Thus only components disjoint from the root and all pure-prime marks may change.  Conversely, a selection of precisely these components satisfies Theorem~\ref{ii:thm:radial}, yielding a radial binary homometric inverse with unchanged root and pure-prime Fourier data.  Lemma~\ref{ii:lem:defect} makes that inverse increasing.  Lemma~\ref{ii:lem:upset} makes it a subgroup union.  Distinctness and exhaustiveness follow from the uniqueness of Fourier inversion.\qedhere
\end{proof}

\begin{remark}[Degenerate cases]\label{ii:rem:degenerate}
The empty set is homometric only to itself, since $C_A(0)=|A|=0$; its graph has no vertices.  For $N=1$ the sole nonempty set is the singleton/full group; the sole conductor is the anchored root.  Zero Fourier coefficients are excluded from the vertex set but were explicitly covered in the adjacent-congruence proof.  Repeated prime factors and all dyadic heights are included in the modular exponents in \eqref{ii:eq:edge}.
\end{remark}

